% Generated from locale/tr/content/computability/computability-theory/total.tex; do not hand-edit.


\olfileid[tr]{cmp}{thy}{tot}
\olsection{Tümellik Karar Verilemezdir}

Bir şeyin hesaplanamaz olduğunu göstermek için $s$-$m$-$n$ teoremini
kullanmanın bir başka örneğini ele alalım. $\fn{Tot}$, tümel hesaplanabilir
fonksiyonların indisleri kümesi olsun:
\[
  \fn{Tot}=\Setabs{x}{\text{her $y$ için }\cfind{x}(y)\fdefined}.
\]

\begin{prop}
\ollabel{prop:total}
$\fn{Tot}$ hesaplanabilir değildir.
\end{prop}

\begin{proof}
$\fn{Tot}$'un hesaplanabilir olmadığını görmek için $K$'nin ona
indirgenebilir olduğunu göstermek yeterlidir. $h(x,y)$ fonksiyonu
\[
  h(x,y)\simeq\begin{cases}
    0 & \text{$x\in K$ ise}\\
    \fundefined & \text{aksi durumda}
  \end{cases}
\]
ile tanımlansın. $h(x,y)$ fonksiyonunun $y$'ye hiç bağlı olmadığına dikkat
edin. $h$'nin kısmi hesaplanabilir olduğunu görmek zor değildir: $x,y$
girdisinde önce $\cfind{x}$ fonksiyonunu $x$ girdisinde benzetiriz; bu hesap
durursa $h(x,y)$, $0$ çıktısını verip durur. Dolayısıyla $\Zero$ sabit sıfır
fonksiyonu olmak üzere $h(x,y)$ yalnızca
$\Zero(\umin{s}{T(x,x,s)})$'dir.

$s$-$m$-$n$ teoremi uyarınca her $x$ ile $y$ için
\[
  \cfind{k(x)}(y)=\begin{cases}
    0 & \text{$x\in K$ ise}\\
    \fundefined & \text{aksi durumda}
  \end{cases}
\]
olacak ilkel özyinelemeli bir $k(x)$ fonksiyonu vardır. Böylece
$\cfind{k(x)}$, $x\in K$ ise tümel, aksi durumda tanımsızdır. Dolayısıyla
$k$, $K$'nin $\fn{Tot}$'a bir indirgemesidir.
\end{proof}

\begin{digress}
Aslında $\fn{Tot}$ hesaplanabilir biçimde numaralandırılabilir bile değildir;
karmaşıklığı ``aritmetik hiyerarşi''de daha yukarıdadır. Fakat burada bu
güçlendirmeyle ilgilenmeyeceğiz.
\end{digress}

